% GENERATED FILE. Edit canonical lessons, then run npm run generate:books.
% canonical-source-hash: fnv1a64:2a0d44558265ef82
% canonical-lessons: KA-C414-kattunittu, KA-C414-garistha, KA-C414-kanistha, KA-C414-nigadipadisu, KA-C414-nirbandhisu

\chapter{Limits}
\label{ch:ka-limits}
\hlchaptermodality{414}

\begin{chapteropening}
\textbf{By the end of this chapter:} \emph{I can say strictness, maximum, minimum, to fix and to restrict.}

Complete the last lesson of chapter 414: I can say strictness, maximum, minimum, to fix and to restrict.
\end{chapteropening}

\section[\texorpdfstring{\kn{ಕಟ್ಟುನಿಟ್ಟು} (kaṭṭuniṭṭu)}{ಕಟ್ಟುನಿಟ್ಟು (kaṭṭuniṭṭu)}]{\kn{ಕಟ್ಟುನಿಟ್ಟು} (kaṭṭuniṭṭu) — strictness}

\label{lesson:KA-C414-kattunittu}



\begin{quote}
Before the new one: say the Kannada for to estimate, then the Kannada for an average.
\end{quote}

\subsection*{You'll want to know: \kn{ಕಟ್ಟುನಿಟ್ಟು}}
\textbf{\kn{ಕಟ್ಟುನಿಟ್ಟು}} — \emph{kaṭṭuniṭṭu} — \textquotedblleft{}strictness\textquotedblright{}.

\begin{quote}
\kn{ಪರೀಕ್ಷೆಯ} \kn{ಸಮಯದಲ್ಲಿ} \kn{ನಿಯಮಗಳು} \kn{ಕಟ್ಟುನಿಟ್ಟಾಗಿ} \kn{ಇರುತ್ತವೆ}.
\end{quote}

\begin{quote}
\emph{parīkṣeya samayadalli niyamagaḷu kaṭṭuniṭṭāgi iruttave.}
\end{quote}

\begin{quote}
During exams the rules are strict.
\end{quote}

\kn{ಕಟ್ಟುನಿಟ್ಟಾದ}, strict; \kn{ಕಟ್ಟುನಿಟ್ಟಾಗಿ}, strictly.

\begin{cousinweb}
\textbf{\kn{ಕಟ್ಟುನಿಟ್ಟು}} (\emph{kaṭṭuniṭṭu}) is native Kannada, a pair-word built on \emph{kaṭṭu}, \textquotedblleft{}to bind, a restraint\textquotedblright{} (DEDR 1147); the second half, \emph{niṭṭu}, is not settled.
\end{cousinweb}

\begin{grammarlens}[title={What you've built}]
One more word for saying what happens if.
\end{grammarlens}

\subsection*{Your turn}
\begin{itemize}
\raggedright
  \item \emph{Say it:} \emph{kaṭṭuniṭṭu}
  \item \emph{Say it:} \emph{kaṭṭuniṭṭu}, once more
  \item \emph{Say it:} \emph{parīkṣeya samayadalli niyamagaḷu kaṭṭuniṭṭāgi iruttave.} — During exams the rules are strict.
  \item \emph{Recall:} say the Kannada for to estimate, then the Kannada for an average, then say \emph{kaṭṭuniṭṭu} again
\end{itemize}

\begin{tcolorbox}[breakable,colback=teal!4,colframe=teal!35!black,title={Before you move on}]
What does \kn{ಕಟ್ಟುನಿಟ್ಟು} mean? (\textquotedblleft{}Strictness\textquotedblright{}.) Say it once more.
\end{tcolorbox}
\section[\texorpdfstring{\kn{ಗರಿಷ್ಠ} (gariṣṭha)}{ಗರಿಷ್ಠ (gariṣṭha)}]{\kn{ಗರಿಷ್ಠ} (gariṣṭha) — maximum}

\label{lesson:KA-C414-garistha}



\begin{quote}
Before the new one: say the Kannada for an average, then the Kannada for strictness.
\end{quote}

\subsection*{You'll want to know: \kn{ಗರಿಷ್ಠ}}
\textbf{\kn{ಗರಿಷ್ಠ}} — \emph{gariṣṭha} — \textquotedblleft{}maximum\textquotedblright{}.

\begin{quote}
\kn{ಈ} \kn{ರಸ್ತೆಯಲ್ಲಿ} \kn{ಗರಿಷ್ಠ} \kn{ವೇಗ} \kn{ನಲವತ್ತು}.
\end{quote}

\begin{quote}
\emph{ī rasteyalli gariṣṭha vēga nalavattu.}
\end{quote}

\begin{quote}
The maximum speed on this road is forty.
\end{quote}

\kn{ಗರಿಷ್ಠ} \kn{ಅಂಕ}, the top mark.

\begin{cousinweb}
\textbf{\kn{ಗರಿಷ್ಠ}} (\emph{gariṣṭha}) is borrowed from Sanskrit \emph{gariṣṭha}, \textquotedblleft{}heaviest\textquotedblright{}, the superlative of \emph{guru}, \textquotedblleft{}heavy\textquotedblright{} (Monier-Williams).
\end{cousinweb}

\begin{grammarlens}[title={What you've built}]
One more word for saying what happens if.
\end{grammarlens}

\subsection*{Your turn}
\begin{itemize}
\raggedright
  \item \emph{Say it:} \emph{gariṣṭha}
  \item \emph{Say it:} \emph{gariṣṭha}, once more
  \item \emph{Say it:} \emph{ī rasteyalli gariṣṭha vēga nalavattu.} — The maximum speed on this road is forty.
  \item \emph{Recall:} say the Kannada for an average, then the Kannada for strictness, then say \emph{gariṣṭha} again
\end{itemize}

\begin{tcolorbox}[breakable,colback=teal!4,colframe=teal!35!black,title={Before you move on}]
What does \kn{ಗರಿಷ್ಠ} mean? (\textquotedblleft{}Maximum\textquotedblright{}.) Say it once more.
\end{tcolorbox}
\section[\texorpdfstring{\kn{ಕನಿಷ್ಠ} (kaniṣṭha)}{ಕನಿಷ್ಠ (kaniṣṭha)}]{\kn{ಕನಿಷ್ಠ} (kaniṣṭha) — minimum; the least}

\label{lesson:KA-C414-kanistha}



\begin{quote}
Before the new one: say the Kannada for strictness, then the Kannada for maximum.
\end{quote}

\subsection*{You'll want to know: \kn{ಕನಿಷ್ಠ}}
\textbf{\kn{ಕನಿಷ್ಠ}} — \emph{kaniṣṭha} — \textquotedblleft{}minimum; the least\textquotedblright{}.

\begin{quote}
\kn{ಈ} \kn{ಕೆಲಸಕ್ಕೆ} \kn{ಕನಿಷ್ಠ} \kn{ಮೂರು} \kn{ದಿನ} \kn{ಬೇಕು}.
\end{quote}

\begin{quote}
\emph{ī kelasakke kaniṣṭha mūru dina bēku.}
\end{quote}

\begin{quote}
This job needs at least three days.
\end{quote}

\kn{ಕನಿಷ್ಠ} \kn{ವೇತನ}, the minimum wage.

\begin{cousinweb}
\textbf{\kn{ಕನಿಷ್ಠ}} (\emph{kaniṣṭha}) is borrowed from Sanskrit \emph{kaniṣṭha}, \textquotedblleft{}the youngest, the least\textquotedblright{} (CDIAL 2718).
\end{cousinweb}

\begin{grammarlens}[title={What you've built}]
One more word for saying what happens if.
\end{grammarlens}

\subsection*{Your turn}
\begin{itemize}
\raggedright
  \item \emph{Say it:} \emph{kaniṣṭha}
  \item \emph{Say it:} \emph{kaniṣṭha}, once more
  \item \emph{Say it:} \emph{ī kelasakke kaniṣṭha mūru dina bēku.} — This job needs at least three days.
  \item \emph{Recall:} say the Kannada for strictness, then the Kannada for maximum, then say \emph{kaniṣṭha} again
\end{itemize}

\begin{tcolorbox}[breakable,colback=teal!4,colframe=teal!35!black,title={Before you move on}]
What does \kn{ಕನಿಷ್ಠ} mean? (\textquotedblleft{}Minimum; the least\textquotedblright{}.) Say it once more.
\end{tcolorbox}
\section[\texorpdfstring{\kn{ನಿಗದಿಪಡಿಸು} (nigadipaḍisu)}{ನಿಗದಿಪಡಿಸು (nigadipaḍisu)}]{\kn{ನಿಗದಿಪಡಿಸು} (nigadipaḍisu) — to fix; to set}

\label{lesson:KA-C414-nigadipadisu}



\begin{quote}
Before the new one: say the Kannada for maximum, then the Kannada for minimum.
\end{quote}

\subsection*{You'll want to know: \kn{ನಿಗದಿಪಡಿಸು}}
\textbf{\kn{ನಿಗದಿಪಡಿಸು}} — \emph{nigadipaḍisu} — \textquotedblleft{}to fix; to set\textquotedblright{}.

\begin{quote}
\kn{ನಾವು} \kn{ಮದುವೆಯ} \kn{ದಿನಾಂಕವನ್ನು} \kn{ನಿಗದಿಪಡಿಸಿದೆವು}.
\end{quote}

\begin{quote}
\emph{nāvu maduveya dināṅkavannu nigadipaḍisidevu.}
\end{quote}

\begin{quote}
We fixed the date of the wedding.
\end{quote}

Past (I …): \kn{ನಿಗದಿಪಡಿಸಿದೆ}, \emph{nigadipaḍiside}.

\kn{ನಿಗದಿತ}, fixed: \kn{ನಿಗದಿತ} \kn{ಸಮಯ}, the set time.

\begin{cousinweb}
\textbf{\kn{ನಿಗದಿಪಡಿಸು}} (\emph{nigadipaḍisu}) is built from \emph{nigadi}, \textquotedblleft{}fixed\textquotedblright{}, with -\emph{paḍisu}, which makes a verb; the origin of \emph{nigadi} is not settled.
\end{cousinweb}

\begin{grammarlens}[title={What you've built}]
One more word for saying what happens if.
\end{grammarlens}

\subsection*{Your turn}
\begin{itemize}
\raggedright
  \item \emph{Say it:} \emph{nigadipaḍisu}
  \item \emph{Say it:} \emph{nigadipaḍisu}, once more
  \item \emph{Say it:} \emph{nāvu maduveya dināṅkavannu nigadipaḍisidevu.} — We fixed the date of the wedding.
  \item \emph{Recall:} say the Kannada for maximum, then the Kannada for minimum, then say \emph{nigadipaḍisu} again
\end{itemize}

\begin{tcolorbox}[breakable,colback=teal!4,colframe=teal!35!black,title={Before you move on}]
What does \kn{ನಿಗದಿಪಡಿಸು} mean? (\textquotedblleft{}To fix; to set\textquotedblright{}.) Say it once more.
\end{tcolorbox}
\section[\texorpdfstring{\kn{ನಿರ್ಬಂಧಿಸು} (nirbandhisu)}{ನಿರ್ಬಂಧಿಸು (nirbandhisu)}]{\kn{ನಿರ್ಬಂಧಿಸು} (nirbandhisu) — to restrict}

\label{lesson:KA-C414-nirbandhisu}



\begin{quote}
Before the new one: say the Kannada for minimum, then the Kannada for to fix.
\end{quote}

\subsection*{You'll want to know: \kn{ನಿರ್ಬಂಧಿಸು}}
\textbf{\kn{ನಿರ್ಬಂಧಿಸು}} — \emph{nirbandhisu} — \textquotedblleft{}to restrict\textquotedblright{}.

\begin{quote}
\kn{ಸರ್ಕಾರ} \kn{ಪ್ಲಾಸ್ಟಿಕ್} \kn{ಬಳಕೆಯನ್ನು} \kn{ನಿರ್ಬಂಧಿಸಿತು}.
\end{quote}

\begin{quote}
\emph{sarkāra plāsṭik baḷakeyannu nirbandhisitu.}
\end{quote}

\begin{quote}
The government restricted the use of plastic.
\end{quote}

Past (I …): \kn{ನಿರ್ಬಂಧಿಸಿದೆ}, \emph{nirbandhiside}.

\kn{ನಿರ್ಬಂಧ}, a restriction.

\begin{cousinweb}
\textbf{\kn{ನಿರ್ಬಂಧಿಸು}} (\emph{nirbandhisu}) is built on Sanskrit \emph{nirbandha}, \textquotedblleft{}perseverance, insistence\textquotedblright{} (CDIAL 7355), from \emph{bandh}, \textquotedblleft{}to bind\textquotedblright{}, with the Kannada verb ending -\emph{isu}.
\end{cousinweb}

\begin{grammarlens}[title={What you've built}]
One more word for saying what happens if.
\end{grammarlens}

\subsection*{Your turn}
\begin{itemize}
\raggedright
  \item \emph{Say it:} \emph{nirbandhisu}
  \item \emph{Say it:} \emph{nirbandhisu}, once more
  \item \emph{Say it:} \emph{sarkāra plāsṭik baḷakeyannu nirbandhisitu.} — The government restricted the use of plastic.
  \item \emph{Recall:} say the Kannada for minimum, then the Kannada for to fix, then say \emph{nirbandhisu} again
\end{itemize}

\begin{tcolorbox}[breakable,colback=teal!4,colframe=teal!35!black,title={Before you move on}]
What does \kn{ನಿರ್ಬಂಧಿಸು} mean? (\textquotedblleft{}To restrict\textquotedblright{}.) Say it once more.
\end{tcolorbox}
